\documentclass[11pt]{article}
\usepackage[a4paper,margin=27mm]{geometry}
\usepackage{amsmath,amssymb,amsthm,xcolor,booktabs,array}
\newcommand{\PROVED}{\textcolor{green!40!black}{\textbf{PROVED}}}
\newcommand{\OPEN}{\textcolor{red!70!black}{\textbf{OPEN}}}
\newcommand{\AUDIT}{\textcolor{orange!70!black}{\textbf{AUDIT}}}
\newtheorem{theorem}{Theorem}
\newtheorem{corollary}{Corollary}
\title{The exact moving-Tate terminal cell (v117)}
\date{13 August 2026}

\begin{document}
\maketitle

\section{An exact Hilbert decomposition of a moving primitive space}

Let $H_0,E$ be Hilbert spaces, let $B_0:H_0\to\mathbb C^r$ be bounded
and onto, and let $B_E:E\to\mathbb C^r$ be bounded.  Put
\begin{equation}
 K_0=\ker B_0,
 \qquad G_0=B_0B_0^*,
 \qquad C=B_0^*G_0^{-1}B_E,
 \qquad M=I_E+C^*C=I_E+B_E^*G_0^{-1}B_E.                  \tag{1}
\end{equation}
The enlarged moment map is
\begin{equation}
 B:H_0\oplus E\longrightarrow\mathbb C^r,
 \qquad B(h,e)=B_0h+B_Ee.                                  \tag{2}
\end{equation}

\begin{theorem}[Moving-constraint terminal cell --- \PROVED]
The enlarged primitive space has the orthogonal decomposition
\begin{equation}
 \ker B=(K_0\oplus0)\ \widehat\oplus\
          \{(-Ce,e):e\in E\}.                              \tag{3}
\end{equation}
Moreover
\begin{equation}
 W:E\longrightarrow H_0\oplus E,
 \qquad Wx=\binom{-CM^{-1/2}x}{M^{-1/2}x}                  \tag{4}
\end{equation}
is an isometry onto the second summand in (3).  Hence, if $P_0$ denotes
the projection onto $K_0$ and $P$ the projection onto $\ker B$, then
\begin{equation}
                         P=(P_0\oplus0)+WW^*.                \tag{5}
\end{equation}
No inverse of an operator compressed to a primitive space occurs in
(1)--(5); the only inverse is the $r\times r$ Gram matrix $G_0$.
\end{theorem}

\begin{proof}
Since $B_0$ is onto,
\[
 H_0=K_0\widehat\oplus\operatorname{Ran}B_0^*.
\]
Write $h=k+r$ according to this decomposition.  The equation $B(h,e)=0$
is $B_0r=-B_Ee$.  The restriction of $B_0$ to
$\operatorname{Ran}B_0^*$ is an isomorphism onto $\mathbb C^r$, and its
inverse is $B_0^*G_0^{-1}$.  Thus $r=-Ce$, proving (3).  The two summands
are orthogonal because $K_0\perp\operatorname{Ran}B_0^*$.

For $x\in E$,
\[
 \|Wx\|^2
 =\langle M^{-1/2}x,(I+C^*C)M^{-1/2}x\rangle
 =\|x\|^2.
\]
The range of $W$ is the graph in (3), so the two orthogonal range
projections add to (5).
\end{proof}

The terminal column (4) satisfies the exact conservation law
\begin{equation}
 \|M^{-1/2}x\|^2+\|CM^{-1/2}x\|^2=\|x\|^2.                 \tag{6}
\end{equation}
In particular, both its newborn port $M^{-1/2}$ and its old Tate return
$CM^{-1/2}$ are contractions.  The column can therefore be completed to
a unitary Julia colligation without imposing any sign on the physical
form.

\section{Application to the two Tate moments}

For $I_T=(-T,T)$ define
\begin{equation}
 B_Tf=inom{\langle f,e^{t/2}\rangle}{
                  \langle f,e^{-t/2}\rangle}.               \tag{7}
\end{equation}
Its Gram matrix is
\begin{equation}
 B_TB_T^*=2
 \begin{pmatrix}\sinh T&T\\T&\sinh T\end{pmatrix},         \tag{8}
\end{equation}
whose eigenvalues $2(\sinh T\pm T)$ are strictly positive for $T>0$.
Let $T'<T+\infty$, identify
\[
 L^2(I_{T'})=L^2(I_T)\oplus L^2(I_{T'}\setminus I_T),
\]
and let $B_E$ be (7) restricted to the newborn set
$E=I_{T'}\setminus I_T$.

\begin{corollary}[Exact Tate threshold update --- \PROVED]
For every $0<T<T'$, the primitive Hilbert space satisfies
\begin{equation}
 \ker B_{T'}=ker B_T\ \widehat\oplus\ W_{T,T'}L^2(E),      \tag{9}
\end{equation}
where $W_{T,T'}$ is given explicitly by (1), (4), and (8).  Thus the two
Tate constraints do not create an uncontrolled moving-projection error:
their entire old/new coupling is the conservative rank-two terminal cell
(4)--(6).
\end{corollary}

\begin{proof}
Apply Theorem 1 with $r=2$, $H_0=L^2(I_T)$, and with the maps (7).
Invertibility follows from (8).
\end{proof}

At the arithmetic thresholds
\begin{equation}
 T=T_N=\tfrac12\log N,
 \qquad T'=T_{N+1},
 \qquad \delta_N=T'-T,                                      \tag{9a}
\end{equation}
the two moment Grams commute.  Indeed,
\begin{equation}
 B_EB_E^*=2
 \begin{pmatrix}
  \sinh T'-\sinh T&\delta_N\\
  \delta_N&\sinh T'-\sinh T
 \end{pmatrix}.                                             \tag{9b}
\end{equation}
The nonzero eigenvalues of $C^*C$ are therefore
\begin{equation}
 \kappa_{N,\pm}=
 \frac{\sinh T'-\sinh T\ \pm\delta_N}
      {\sinh T\ \pm T}.                                    \tag{9c}
\end{equation}

\begin{corollary}[Vanishing Tate return on thin arithmetic cells ---
\PROVED]
As $N\to\infty$,
\begin{equation}
 \|C_{T_N,T_{N+1}}\|^2
 =\max(\kappa_{N,+},\kappa_{N,-})
 =\frac1{2N}+o(N^{-1}).                                     \tag{9d}
\end{equation}
Consequently
\begin{equation}
 \|CM^{-1/2}\|=O(N^{-1/2}),
 \qquad \|M^{-1/2}-I\|=O(N^{-1}).                           \tag{9e}
\end{equation}
Thus the moving Tate terminal is not an unbounded threshold connection;
its old return vanishes at the rate $N^{-1/2}$.
\end{corollary}

\begin{proof}
The nonzero spectra of $B_E^*G_0^{-1}B_E$ and
$G_0^{-1/2}B_EB_E^*G_0^{-1/2}$ agree.  Equations (8)--(9b) diagonalize
both matrices in the vectors $(1,1)$ and $(1,-1)$, giving (9c).  Since
$\delta_N=\frac12\log(1+N^{-1})=\frac1{2N}+O(N^{-2})$ and
\[
 \sinh(T+\delta_N)-\sinh T
 =\delta_N\cosh T+O(\delta_N^2\sinh T),
\]
while $(\cosh T\pm1)/(\sinh T\pm T)\to1$, equation (9d) follows.
Finally $M=I+C^*C$ and functional calculus give (9e).
\end{proof}

\section{Exact compression of an arbitrary physical block}

Let a bounded self-adjoint operator on $H_0\oplus E$ be written
\begin{equation}
 \mathcal A=
 \begin{pmatrix}A_{00}&A_{0E}\\A_{E0}&A_{EE}\end{pmatrix},
 \qquad N=M^{-1/2}.                                         \tag{10}
\end{equation}
Under the unitary identification
$K_0\oplus E\to\ker B$, $(k,x)\mapsto(k,0)+Wx$, its primitive
compression is exactly
\begin{equation}
 \begin{pmatrix}
  P_0A_{00}P_0& P_0(A_{0E}-A_{00}C)N\\
  N(A_{E0}-C^*A_{00})P_0&
  N(A_{EE}-C^*A_{0E}-A_{E0}C+C^*A_{00}C)N
 \end{pmatrix}.                                             \tag{11}
\end{equation}

\begin{theorem}[No hidden moving-Tate correction --- \PROVED]
Formula (11) is the complete old/new block matrix after primitiveness.
In particular, every correction caused by the moving Tate projection is
one of the four displayed terms containing $C$; there is no additional
connection term.
\end{theorem}

\begin{proof}
Let
\[
 V=\begin{pmatrix}P_0&-CN\\0&N\end{pmatrix}:K_0\oplus E
        \longrightarrow H_0\oplus E.
\]
By Theorem 1, $V$ is unitary onto $\ker B$.  Direct multiplication of
$V^*\mathcal AV$ gives (11).
\end{proof}

Formula (11) isolates the next, genuinely physical calculation.  For the
Gamma--Euler--Tate operator $\mathcal A_{T'}$, one must factor the Schur
innovation
\begin{align}
 \mathcal S_{T,T'}={}&
 N(A_{EE}-C^*A_{0E}-A_{E0}C+C^*A_{00}C)N \notag\\
 &-N(A_{E0}-C^*A_{00})P_0
   (P_0A_{00}P_0)^\dagger
   P_0(A_{0E}-A_{00}C)N                                    \tag{12}
\end{align}
on the support on which the Moore--Penrose expression is defined.
The Tate part of (12) is now explicit and conservative by (1), (6), and
(8); it may not be cited as the source of an unspecified residual.

\begin{center}
\begin{tabular}{>{\raggedright\arraybackslash}p{.55\linewidth}
                >{\centering\arraybackslash}p{.15\linewidth}
                >{\raggedright\arraybackslash}p{.20\linewidth}}
\toprule
Statement & Status & Location \\
\midrule
Orthogonal old-primitive/new-terminal decomposition & \PROVED & (3)--(5) \\
Exact Hilbert conservation of the Tate terminal & \PROVED & (6) \\
Explicit two-moment threshold cell & \PROVED & (7)--(9) \\
Vanishing old Tate return on arithmetic cells & \PROVED & (9a)--(9e) \\
Complete moving-primitive compression formula & \PROVED & (11) \\
Positive factorization of the physical innovation (12) & \OPEN & Euler--Gamma calculation \\
Condition $G$, row (d) & \OPEN & follows only after (12) for all cells \\
\bottomrule
\end{tabular}
\end{center}

\end{document}
